\scanpage{23}
On généralise la formule fondamentale au cas d'un \uline{morphisme propre quelconque}
\diagram{\begin{tikzcd}[ampersand replacement=\&]f:\quad X \arrow[r] \& Y\end{tikzcd}}
deux variétés \deleted{\unsure{algébr.}} holomorphes, \unsure{partout|partant} de codim relative $p$ :
\diagram{\fbox{$\begin{tikzcd}[ampersand replacement=\&,row sep=large,column sep=huge]
 H^{\unsure{2i}}(X,\mathbf C) \arrow[r,"f_*^{\mathrm{top}}"] \arrow[d,dash,"\wr"] \& H^{\unsure{2i+2p}}(Y,\mathbf C) \arrow[d,dash,"\wr"] \\
 H^{\unsure{2i}}(X,\Omega_X^\bullet) \arrow[r,"{\frac{1}{(2i\pi)^p}f_*^{\mathrm{an}}}"] \& H^{\unsure{2i+2p}}(Y,\Omega_Y^\bullet)
\end{tikzcd}$}}

D'où généralisation analogue pour des variétés \unsure{algébriques} sur $\mathbf C$ ....

\mspage{22}
